% !TeX encoding = UTF-8
% !TeX program = xelatex
% !TeX spellcheck = en_US

\documentclass[degree=master]{thuthesis}
  % 学位 degree:
  %   doctor | master | bachelor | postdoc
  % 学位类型 degree-type:
  %   academic（默认）| professional
  % 语言 language
  %   chinese（默认）| english
  % 字体库 fontset
  %   windows | mac | fandol | ubuntu
  % 建议终版使用 Windows 平台的字体编译


% 论文基本配置，加载宏包等全局配置
% !TeX root = ./thuthesis-example.tex

% 论文基本信息配置

\thusetup{
  %******************************
  % 注意：
  %   1. 配置里面不要出现空行
  %   2. 不需要的配置信息可以删除
  %   3. 建议先阅读文档中所有关于选项的说明
  %******************************
  %
  % 输出格式
  %   选择打印版（print）或用于提交的电子版（electronic），前者会插入空白页以便直接双面打印
  %
  output = print,
  % 格式类型
  %   默认为论文（thesis），也可以设置为开题报告（proposal）
  % thesis-type = proposal,
  %
  % 标题
  %   可使用“\\”命令手动控制换行
  %
  title  = {清华大学学位论文 \LaTeX{} 模板\\使用示例文档 v\version},
  title* = {An Introduction to \LaTeX{} Thesis Template of Tsinghua
            University v\version},
  %
  % 学科门类
  %   1. 学术型
  %      - 中文
  %        需注明所属的学科门类，例如：
  %        哲学、经济学、法学、教育学、文学、历史学、理学、工学、农学、医学、
  %        军事学、管理学、艺术学
  %      - 英文
  %        博士：Doctor of Philosophy
  %        硕士：
  %          哲学、文学、历史学、法学、教育学、艺术学门类，公共管理学科
  %          填写“Master of Arts“，其它填写“Master of Science”
  %   2. 专业型
  %      直接填写专业学位的名称，例如：
  %      教育博士、工程硕士等
  %      Doctor of Education, Master of Engineering
  %   3. 本科生不需要填写
  %
  degree-category  = {工学硕士},
  degree-category* = {Master of Science},
  %
  % 培养单位
  %   填写所属院系的全名
  %
  department = {计算机科学与技术系},
  %
  % 学科
  %   1. 研究生学术型学位，获得一级学科授权的学科填写一级学科名称，其他填写二级学科名称
  %   2. 本科生填写专业名称，第二学位论文需标注“（第二学位）”
  %
  discipline  = {计算机科学与技术},
  discipline* = {Computer Science and Technology},
  %
  % 专业领域
  %   1. 设置专业领域的专业学位类别，填写相应专业领域名称
  %   2. 2019 级及之前工程硕士学位论文，在 `engineering-field` 填写相应工程领域名称
  %   3. 其他专业学位类别的学位论文无需此信息
  %
  % professional-field  = {计算机技术},
  % professional-field* = {Computer Technology},
  %
  % 姓名
  %
  author  = {薛瑞尼},
  author* = {Xue Ruini},
  %
  % 学号
  % 仅当书写开题报告时需要（同时设置 `thesis-type = proposal'）
  %
  % student-id = {2000310000},
  %
  % 指导教师
  %   中文姓名和职称之间以英文逗号“,”分开，下同
  %
  supervisor  = {郑纬民, 教授},
  supervisor* = {Professor Zheng Weimin},
  %
  % 副指导教师
  %
  associate-supervisor  = {陈文光, 教授},
  associate-supervisor* = {Professor Chen Wenguang},
  %
  % 联合指导教师
  %
  % co-supervisor  = {某某某, 教授},
  % co-supervisor* = {Professor Mou Moumou},
  %
  % 日期
  %   使用 ISO 格式；默认为当前时间
  %
  % date = {2019-07-07},
  %
  % 是否在中文封面后的空白页生成书脊（默认 false）
  %
  include-spine = false,
  %
  % 密级和年限
  %   秘密, 机密, 绝密
  %
  % secret-level = {秘密},
  % secret-year  = {10},
  %
  % 博士后专有部分
  %
  % clc                = {分类号},
  % udc                = {UDC},
  % id                 = {编号},
  % discipline-level-1 = {计算机科学与技术},  % 流动站（一级学科）名称
  % discipline-level-2 = {系统结构},          % 专业（二级学科）名称
  % start-date         = {2011-07-01},        % 研究工作起始时间
}

% 载入所需的宏包

% 定理类环境宏包
\usepackage{amsthm}
% 也可以使用 ntheorem
% \usepackage[amsmath,thmmarks,hyperref]{ntheorem}

\thusetup{
  %
  % 数学字体
  % math-style = GB,  % GB | ISO | TeX
  math-font  = xits,  % stix | xits | libertinus
}

% 可以使用 nomencl 生成符号和缩略语说明
% \usepackage{nomencl}
% \makenomenclature

% 表格加脚注
\usepackage{threeparttable}

% 表格中支持跨行
\usepackage{multirow}

% 固定宽度的表格。
% \usepackage{tabularx}

% 跨页表格
\usepackage{longtable}

% 算法
\usepackage{algorithm}
\usepackage{algorithmic}

% 量和单位
\usepackage{siunitx}

% 参考文献使用 BibTeX + natbib 宏包
% 顺序编码制
\usepackage[sort]{natbib}
\bibliographystyle{thuthesis-numeric}

% 著者-出版年制
% \usepackage{natbib}
% \bibliographystyle{thuthesis-author-year}

% 生命科学学院要求使用 Cell 参考文献格式（2023 年以前使用 author-date 格式）
% \usepackage{natbib}
% \bibliographystyle{cell}

% 本科生参考文献的著录格式
% \usepackage[sort]{natbib}
% \bibliographystyle{thuthesis-bachelor}

% 参考文献使用 BibLaTeX 宏包
% \usepackage[style=thuthesis-numeric]{biblatex}
% \usepackage[style=thuthesis-author-year]{biblatex}
% \usepackage[style=gb7714-2015]{biblatex}
% \usepackage[style=apa]{biblatex}
% \usepackage[style=mla-new]{biblatex}
% 声明 BibLaTeX 的数据库
% \addbibresource{ref/refs.bib}

% 定义所有的图片文件在 figures 子目录下
\graphicspath{{figures/}}

% 数学命令
\makeatletter
\newcommand\dif{%  % 微分符号
  \mathop{}\!%
  \ifthu@math@style@TeX
    d%
  \else
    \mathrm{d}%
  \fi
}
\makeatother

% hyperref 宏包在最后调用
\usepackage{hyperref}


%----- 定义页眉-----
\makeatletter
\newcommand\@header{}
\newcommand\header[1]{\def\@header{#1}}
\makeatother

%----- 页眉页脚 -----
\usepackage{fancyhdr}
\makeatletter
\pagestyle{fancy}
\fancyhf{}
\fancyhead[C]{\small\kaishu\@header}
\fancyfoot[C]{\small\thepage}
% \renewcommand{\headrulewidth}{0.5pt}
% % \renewcommand{\headrulewidth}{0.5pt}
% \renewcommand{\headrulewidth}{2pt}
\newcommand{\makeheadrule}{%
\makebox[0pt][l]{\rule[0.6\baselineskip]{\headwidth}{1.5pt}}%
\rule[0.7\baselineskip]{\headwidth}{0.5pt}}
\renewcommand{\headrule}{%
{\if@fancyplain\let\headrulewidth\plainheadrulewidth\fi
\makeheadrule}}
\fancyhead[C]{哈尔滨工程大学本科生课程设计}
\makeatother
\begin{document}

% % 封面
% \maketitle

% % 学位论文指导小组、公开评阅人和答辩委员会名单
% % 本科生不需要
% % !TeX root = ../thuthesis-example.tex

\begin{committee}[name={学位论文指导小组、公开评阅人和答辩委员会名单}]

  \newcolumntype{C}[1]{@{}>{\centering\arraybackslash}p{#1}}

  \section*{指导小组名单}

  \begin{center}
    \begin{tabular}{C{3cm}C{3cm}C{9cm}@{}}
      李XX & 教授     & 清华大学 \\
      王XX & 副教授   & 清华大学 \\
      张XX & 助理教授 & 清华大学 \\
    \end{tabular}
  \end{center}


  \section*{公开评阅人名单}

  \begin{center}
    \begin{tabular}{C{3cm}C{3cm}C{9cm}@{}}
      刘XX & 教授   & 清华大学                    \\
      陈XX & 副教授 & XXXX大学                    \\
      杨XX & 研究员 & 中国XXXX科学院XXXXXXX研究所 \\
    \end{tabular}
  \end{center}


  \section*{答辩委员会名单}

  \begin{center}
    \begin{tabular}{C{2.75cm}C{2.98cm}C{4.63cm}C{4.63cm}@{}}
      主席 & 赵XX                  & 教授                    & 清华大学       \\
      委员 & 刘XX                  & 教授                    & 清华大学       \\
          & \multirow{2}{*}{杨XX} & \multirow{2}{*}{研究员} & 中国XXXX科学院 \\
          &                       &                         & XXXXXXX研究所  \\
          & 黄XX                  & 教授                    & XXXX大学       \\
          & 周XX                  & 副教授                  & XXXX大学       \\
      秘书 & 吴XX                  & 助理研究员              & 清华大学       \\
    \end{tabular}
  \end{center}

\end{committee}



% 也可以导入 Word 版转的 PDF 文件
% \begin{committee}[file=figures/committee.pdf]
% \end{committee}


% % 使用授权的说明
% % 本科生开题报告不需要
% \copyrightpage
% % 将签字扫描后授权文件 scan-copyright.pdf 替换原始页面
% % \copyrightpage[file=scan-copyright.pdf]

\frontmatter
% % !TeX root = ../thuthesis-example.tex

% % 中英文摘要和关键字

% \begin{abstract}
%   论文的摘要是对论文研究内容和成果的高度概括。
%   摘要应对论文所研究的问题及其研究目的进行描述，对研究方法和过程进行简单介绍，对研究成果和所得结论进行概括。
%   摘要应具有独立性和自明性，其内容应包含与论文全文同等量的主要信息。
%   使读者即使不阅读全文，通过摘要就能了解论文的总体内容和主要成果。

%   论文摘要的书写应力求精确、简明。
%   切忌写成对论文书写内容进行提要的形式，尤其要避免“第 1 章……；第 2 章……；……”这种或类似的陈述方式。

%   关键词是为了文献标引工作、用以表示全文主要内容信息的单词或术语。
%   关键词不超过 5 个，每个关键词中间用分号分隔。

%   % 关键词用“英文逗号”分隔，输出时会自动处理为正确的分隔符
%   \thusetup{
%     keywords = {关键词 1, 关键词 2, 关键词 3, 关键词 4, 关键词 5},
%   }
% \end{abstract}

% \begin{abstract*}
%   An abstract of a dissertation is a summary and extraction of research work and contributions.
%   Included in an abstract should be description of research topic and research objective, brief introduction to methodology and research process, and summary of conclusion and contributions of the research.
%   An abstract should be characterized by independence and clarity and carry identical information with the dissertation.
%   It should be such that the general idea and major contributions of the dissertation are conveyed without reading the dissertation.

%   An abstract should be concise and to the point.
%   It is a misunderstanding to make an abstract an outline of the dissertation and words “the first chapter”, “the second chapter” and the like should be avoided in the abstract.

%   Keywords are terms used in a dissertation for indexing, reflecting core information of the dissertation.
%   An abstract may contain a maximum of 5 keywords, with semi-colons used in between to separate one another.

%   % Use comma as separator when inputting
%   \thusetup{
%     keywords* = {keyword 1, keyword 2, keyword 3, keyword 4, keyword 5},
%   }
% \end{abstract*}


% 目录
\tableofcontents

% % 插图和附表清单
% % 本科生的插图索引和表格索引需要移至正文之后、参考文献前
% % \listoffiguresandtables  % 插图和附表清单（仅限研究生）
% \listoffigures           % 插图清单
% \listoftables            % 附表清单

% 符号对照表
% !TeX root = ../thuthesis-example.tex

\begin{denotation}[3cm]
  \item[PI] 聚酰亚胺
  \item[MPI] 聚酰亚胺模型化合物，N-苯基邻苯酰亚胺
  \item[PBI] 聚苯并咪唑
  \item[MPBI] 聚苯并咪唑模型化合物，N-苯基苯并咪唑
  \item[PY] 聚吡咙
  \item[PMDA-BDA] 均苯四酸二酐与联苯四胺合成的聚吡咙薄膜
  \item[MPY] 聚吡咙模型化合物
  \item[As-PPT] 聚苯基不对称三嗪
  \item[MAsPPT] 聚苯基不对称三嗪单模型化合物，3,5,6-三苯基-1,2,4-三嗪
  \item[DMAsPPT] 聚苯基不对称三嗪双模型化合物（水解实验模型化合物）
  \item[S-PPT] 聚苯基对称三嗪
  \item[MSPPT] 聚苯基对称三嗪模型化合物，2,4,6-三苯基-1,3,5-三嗪
  \item[PPQ] 聚苯基喹噁啉
  \item[MPPQ] 聚苯基喹噁啉模型化合物，3,4-二苯基苯并二嗪
  \item[HMPI] 聚酰亚胺模型化合物的质子化产物
  \item[HMPY] 聚吡咙模型化合物的质子化产物
  \item[HMPBI] 聚苯并咪唑模型化合物的质子化产物
  \item[HMAsPPT] 聚苯基不对称三嗪模型化合物的质子化产物
  \item[HMSPPT] 聚苯基对称三嗪模型化合物的质子化产物
  \item[HMPPQ] 聚苯基喹噁啉模型化合物的质子化产物
  \item[PDT] 热分解温度
  \item[HPLC] 高效液相色谱（High Performance Liquid Chromatography）
  \item[HPCE] 高效毛细管电泳色谱（High Performance Capillary lectrophoresis）
  \item[LC-MS] 液相色谱-质谱联用（Liquid chromatography-Mass Spectrum）
  \item[TIC] 总离子浓度（Total Ion Content）
  \item[\textit{ab initio}] 基于第一原理的量子化学计算方法，常称从头算法
  \item[DFT] 密度泛函理论（Density Functional Theory）
  \item[$E_a$] 化学反应的活化能（Activation Energy）
  \item[ZPE] 零点振动能（Zero Vibration Energy）
  \item[PES] 势能面（Potential Energy Surface）
  \item[TS] 过渡态（Transition State）
  \item[TST] 过渡态理论（Transition State Theory）
  \item[$\increment G^\neq$] 活化自由能（Activation Free Energy）
  \item[$\kappa$] 传输系数（Transmission Coefficient）
  \item[IRC] 内禀反应坐标（Intrinsic Reaction Coordinates）
  \item[$\nu_i$] 虚频（Imaginary Frequency）
  \item[ONIOM] 分层算法（Our own N-layered Integrated molecular Orbital and molecular Mechanics）
  \item[SCF] 自洽场（Self-Consistent Field）
  \item[SCRF] 自洽反应场（Self-Consistent Reaction Field）
\end{denotation}



% 也可以使用 nomencl 宏包，需要在导言区
% \usepackage{nomencl}
% \makenomenclature

% 在这里输出符号说明
% \printnomenclature[3cm]

% 在正文中的任意为都可以标题
% \nomenclature{PI}{聚酰亚胺}
% \nomenclature{MPI}{聚酰亚胺模型化合物，N-苯基邻苯酰亚胺}
% \nomenclature{PBI}{聚苯并咪唑}
% \nomenclature{MPBI}{聚苯并咪唑模型化合物，N-苯基苯并咪唑}
% \nomenclature{PY}{聚吡咙}
% \nomenclature{PMDA-BDA}{均苯四酸二酐与联苯四胺合成的聚吡咙薄膜}
% \nomenclature{MPY}{聚吡咙模型化合物}
% \nomenclature{As-PPT}{聚苯基不对称三嗪}
% \nomenclature{MAsPPT}{聚苯基不对称三嗪单模型化合物，3,5,6-三苯基-1,2,4-三嗪}
% \nomenclature{DMAsPPT}{聚苯基不对称三嗪双模型化合物（水解实验模型化合物）}
% \nomenclature{S-PPT}{聚苯基对称三嗪}
% \nomenclature{MSPPT}{聚苯基对称三嗪模型化合物，2,4,6-三苯基-1,3,5-三嗪}
% \nomenclature{PPQ}{聚苯基喹噁啉}
% \nomenclature{MPPQ}{聚苯基喹噁啉模型化合物，3,4-二苯基苯并二嗪}
% \nomenclature{HMPI}{聚酰亚胺模型化合物的质子化产物}
% \nomenclature{HMPY}{聚吡咙模型化合物的质子化产物}
% \nomenclature{HMPBI}{聚苯并咪唑模型化合物的质子化产物}
% \nomenclature{HMAsPPT}{聚苯基不对称三嗪模型化合物的质子化产物}
% \nomenclature{HMSPPT}{聚苯基对称三嗪模型化合物的质子化产物}
% \nomenclature{HMPPQ}{聚苯基喹噁啉模型化合物的质子化产物}
% \nomenclature{PDT}{热分解温度}
% \nomenclature{HPLC}{高效液相色谱（High Performance Liquid Chromatography）}
% \nomenclature{HPCE}{高效毛细管电泳色谱（High Performance Capillary lectrophoresis）}
% \nomenclature{LC-MS}{液相色谱-质谱联用（Liquid chromatography-Mass Spectrum）}
% \nomenclature{TIC}{总离子浓度（Total Ion Content）}
% \nomenclature{\textit{ab initio}}{基于第一原理的量子化学计算方法，常称从头算法}
% \nomenclature{DFT}{密度泛函理论（Density Functional Theory）}
% \nomenclature{$E_a$}{化学反应的活化能（Activation Energy）}
% \nomenclature{ZPE}{零点振动能（Zero Vibration Energy）}
% \nomenclature{PES}{势能面（Potential Energy Surface）}
% \nomenclature{TS}{过渡态（Transition State）}
% \nomenclature{TST}{过渡态理论（Transition State Theory）}
% \nomenclature{$\increment G^\neq$}{活化自由能（Activation Free Energy）}
% \nomenclature{$\kappa$}{传输系数（Transmission Coefficient）}
% \nomenclature{IRC}{内禀反应坐标（Intrinsic Reaction Coordinates）}
% \nomenclature{$\nu_i$}{虚频（Imaginary Frequency）}
% \nomenclature{ONIOM}{分层算法（Our own N-layered Integrated molecular Orbital and molecular Mechanics）}
% \nomenclature{SCF}{自洽场（Self-Consistent Field）}
% \nomenclature{SCRF}{自洽反应场（Self-Consistent Reaction Field）}



% 正文部分
\mainmatter
% !TeX root = ../thuthesis-example.tex
\chapter{绪论}
固体火箭发动机是火箭和导弹中广泛应用的动力装置，为了实现我国国防现代化建设大业，作为相关专业的学生应该具备设计一款性能符合要求结构安全可靠的固体火箭发动机的能力。

固体火箭发动机由于具有结构简单、体积小、工作可靠、操作简便、能够长期贮存等优点，即使在今天仍然没有被淘汰依旧具有很大的研究学习价值。
\section{固体火箭发动机的结构}
固体火箭发动机主要由燃烧室、装药、点火装置和喷管等部件组成.
\begin{enumerate}[leftmargin=2em]
  \item 燃烧室：
  
  发动机中的关键装置，提供推进剂燃烧空间，将化学能转化为高温高压的热能，并采用高性能金属或复合材料作为壳体承受高温高压，同时周围涂有绝热层采取隔热措施保护壳体。
  \item 装药：
  
  固体推进剂制成的药柱，作为发动机的能源和工质源，需具备特定几何形状，并通过控制燃烧表面变化，以实现预期推力。包括\textbf{双基(double base，DB)推进剂}和\textbf{复合推进剂(compositepropellant)}。

  \item 点火装置：
  
  由保险机构和点火器组成，用于点燃发动机装药，使发动机顺利起动。点火器中有接受起动信息就开始工作的始发器如电发火管，还有相当数量的点火药。对于大型的固体发动机，其点火器往往设计成了小型的点火发动机。

  \item 喷管：
  
  发动机中关键的能量转换装置，通过加速膨胀高温高压燃气以产生推力，同时也并具备推力向量控制和推力终止装置来实现飞行器的姿态控制和精确停机。
\end{enumerate}
% \section{固体火箭发动机设计总要求}


% 火箭发动机课程设计是飞行器动力工程专业的必修实践课，旨在让学生掌握固体发动机设计的基本方法和原理。课程任务包括：根据给定设计要求选择药形，进行装药设计，计算发动机内弹道性能，分析装药特点，并筛选出最佳方案。

% 选择一种药形，采用附录A中表1-3所提供的推进剂，设计一组装药，进行发动机内弹道计算，通过内弹道分析和装药性能对比，最终确定一组装药。题目为“xx形装药设计”。

% \section{固体火箭发动机设计详细参数要求}

% \subsection{装药几何参数}

% \begin{table}
%   \centering
%   \caption{装药几何参数}
%   \begin{tabular}{ll}
%     \toprule
%     相关几何参数          & 要求                         \\
%     \midrule
%     外径   & 230～300mm \\
%     长度   & 1300~1800mm                     \\
%     总冲   & 不小于230kN·s    \\
%     最大推力   & 38～72kN    \\
%     平均推力   & 35～65kN    \\
%     最小推力   & 31～58kN    \\
%     平均工作时间   & 4～7s    \\
%     工作环境温度   & 20℃    \\

%     \bottomrule
%   \end{tabular}
%   \label{tab:three-line}
% \end{table}

% \subsection{进剂种类选择}

% 在附录A表1~3中选取.

% \subsection{壳体材料选择}

% \thusetup{
%   cite-style = inline,
% }

% 参见参考资料\citep{wangyouyuan}。
\section{固体火箭发动机设计原始参数}
\begin{enumerate}[leftmargin=2em]
  \item 发动机的用途
  
  发动机用在何种火箭或导弹上。
  \item 发动机的总冲
  
  总冲量$I$用发动机推力$F$对时间变量$t$在整个发动机工作时间$t_{a}$区间内的积分来表示，即
  \[
    I=\int_{0}^{t_{a}} Fdt
  \]
  \item 发动机的比冲
  
  单位重量推进剂产生的冲量，即
  \[
    I_{sp}=\frac{I}{m_{p}g} 
  \]
  \item 发动机的使用温度范围
  
  由设计要求部门提供。

  \item 发动机的平均推力
  
  会给出在常温下的平均推力$t_{a}$，有时也会给出所限制的最长工作时间$t_{amax}$与最短工作时间$t_{amin}$。

  \item 发动机的工作时间
  
  会给出在常温下的工作时间$\bar{F} $，有时也会给出所限制的最大推力$F_{max}$与最小推力$F_{min}$。

  \item 发动机的推力方案
  
  发动机选择等推力还是变推力。

  \item 发动机的质量限制
  
  限制发动机的总质量$m$和结构质量$m_{m}$。

  \item 发动机的尺寸限制
  
  对发动机燃烧室直径$D$、长径比$L/D$的限制。

  \item 发动机性能的偏差量
  
  括推力、工作时间和总冲的偏差量，和允许的推力偏角等。

  \item 点火延迟期
  
  点火延迟期$t_{ig}$是以推力（或压力）达到其额定值的75$\sim$80\%所需的时间给出的

  \item 对推力矢量控制装置的要求
  
  如有无推力矢量控制装置，推力矢量控制装置的最大侧向力最大偏转角的大小和频率响应等。

  \item 对推力终止装置的要求
  
  如有无推力终止装置，推力终止的时间和作用效果等。

  \item 发动机的贮存期限
  
  有关部门提出贮存要求。

  \item 发动机的运输条件
  
  包括运输方式、运输距离和行车速度等。

  \item 特殊要求
  

  
\end{enumerate}
\section{固体火箭发动机设计任务要求}

固体火箭发动机设计内容包括：

\begin{itemize}[leftmargin=2em]
  \item 发动机总体设计
  \item 发动机装药设计
  \item 发动机燃烧室设计
  \item 发动机喷管设计
  \item 发动机点火装置设计
  \item 推力矢量控制装置设计
  \item 推力终止装置设计
\end{itemize}

固体火箭发动机设计需要考虑的因素

\begin{itemize}[leftmargin=2em]
  \item 良好的可靠性
  \item 良好的安全性
  \item 经济型
\end{itemize}

% !TeX root = ../thuthesis-example.tex

\chapter{固体火箭发动机课程设计}

装药是指填入燃烧室的推进剂药柱，它的形状和尺寸对发动机的性能至关重要。装药的几何特征直接影响燃气的生成速率和变化规律，从而决定了发动机推力和压力随时间的变化。

\section{课程设计的内容}

\begin{enumerate}[leftmargin=2em]
    \item 装药药形选择
    \item 按给定的原始参数计算装药尺寸（多组）
    \item 给出燃面变化规律（通过曲线给出，肉厚步长取1mm）
    \item 计算对应装药无侵蚀燃烧时的发动机工作$p_{t}$曲线
    \item 计算对应装药有侵蚀燃烧时的发动机工作$p_{t}$曲线
    \item 分析侵蚀燃烧对内弹道的影响
    \item 通过内弹道分析和装药性能分析确定一种装药
\end{enumerate}
\section{固体火箭发动机设计详细参数要求}

\subsection{装药几何参数}

\begin{table}
  \centering
  \caption{装药几何参数}
  \begin{tabular}{ll}
    \toprule
    相关几何参数          & 要求                         \\
    \midrule
    外径   & 230～300mm \\
    长度   & 1300~1800mm                     \\
    总冲   & 不小于230kN·s    \\
    最大推力   & 38～72kN    \\
    平均推力   & 35～65kN    \\
    最小推力   & 31～58kN    \\
    平均工作时间   & 4～7s    \\
    工作环境温度   & 20℃    \\

    \bottomrule
  \end{tabular}
  \label{tab:three-line}
\end{table}

\subsection{进剂种类选择}

在附录A表1~3中选取。

\subsection{壳体材料选择}

\thusetup{
  cite-style = inline,
}

参见参考资料\citep{wangyouyuan}。

\section{侵蚀燃烧的计算}

\subsection{星型装药}

对于星型装药，我们采用æ准则拟合沿装药长度平均侵蚀比：
\[
    \overline{\varepsilon }=\overline{\varepsilon }\left( \varphi \right) =\left\{ \begin{matrix}
        1&		(\varphi \le 72.9)\\
        1.3128-1.3249\times 10^{-2}\varphi +1.5527\times 10^{-4}\varphi ^2-4.3868\times 10^{-7}\varphi ^3&		(\varphi >72.9)\\
    \end{matrix} \right. 
\]

式中：$\varphi$——燃通比，$\varphi=A_{b}/A{p}$

\qquad\quad$A_{b}$——燃面，$m^{2}$

\qquad\quad$A{p}$——通气面积，$m^{2}$

\subsection{多孔装药}

对于多孔装药，我们采用下式拟合沿装药长度平均侵蚀比：
\[
    \overline{\varepsilon }=\overline{\varepsilon }\left( \varphi \right) =\left\{ \begin{matrix}
        1&		(\varphi \le 72.9)\\
        1.3128-1.3249\times 10^{-2}\varphi +1.5527\times 10^{-4}\varphi ^2-4.3868\times 10^{-7}\varphi ^3&		(\varphi >72.9)\\
    \end{matrix} \right. 
\]
% !TeX root = ../thuthesis-example.tex

\chapter{数学符号和公式}

\section{数学符号}

中文论文的数学符号默认遵循 GB/T 3102.11—1993《物理科学和技术中使用的数学符号》
\footnote{原 GB 3102.11—1993，自 2017 年 3 月 23 日起，该标准转为推荐性标准。}。
该标准参照采纳 ISO 31-11:1992 \footnote{目前已更新为 ISO 80000-2:2019。}，
但是与 \TeX{} 默认的美国数学学会（AMS）的符号习惯有所区别。
具体地来说主要有以下差异：
\begin{enumerate}
  \item 大写希腊字母默认为斜体，如
    \begin{equation*}
      \Gamma \Delta \Theta \Lambda \Xi \Pi \Sigma \Upsilon \Phi \Psi \Omega.
    \end{equation*}
    注意有限增量符号 $\increment$ 固定使用正体，模板提供了 \cs{increment} 命令。
  \item 小于等于号和大于等于号使用倾斜的字形 $\le$、$\ge$。
  \item 积分号使用正体，比如 $\int$、$\oint$。
  \item
    偏微分符号 $\partial$ 使用正体。
  \item
    省略号 \cs{dots} 按照中文的习惯固定居中，比如
    \begin{equation*}
      1, 2, \dots, n \quad 1 + 2 + \dots + n.
    \end{equation*}
  \item
    实部 $\Re$ 和虚部 $\Im$ 的字体使用罗马体。
\end{enumerate}

以上数学符号样式的差异可以在模板中统一设置。
另外国标还有一些与 AMS 不同的符号使用习惯，需要用户在写作时进行处理：
\begin{enumerate}
  \item 数学常数和特殊函数名用正体，如
    \begin{equation*}
      \uppi = 3.14\dots; \quad
      \symup{i}^2 = -1; \quad
      \symup{e} = \lim_{n \to \infty} \left( 1 + \frac{1}{n} \right)^n.
    \end{equation*}
  \item 微分号使用正体，比如 $\dif y / \dif x$。
  \item 向量、矩阵和张量用粗斜体（\cs{symbf}），如 $\symbf{x}$、$\symbf{\Sigma}$、$\symbfsf{T}$。
  \item 自然对数用 $\ln x$ 不用 $\log x$。
\end{enumerate}


英文论文的数学符号使用 \TeX{} 默认的样式。
如果有必要，也可以通过设置 \verb|math-style| 选择数学符号样式。

关于量和单位推荐使用
\href{http://mirrors.ctan.org/macros/latex/contrib/siunitx/siunitx.pdf}{\pkg{siunitx}}
宏包，
可以方便地处理希腊字母以及数字与单位之间的空白，
比如：
\SI{6.4e6}{m}，
\SI{9}{\micro\meter}，
\si{kg.m.s^{-1}}，
\SIrange{10}{20}{\degreeCelsius}。



\section{数学公式}

数学公式可以使用 \env{equation} 和 \env{equation*} 环境。
注意数学公式的引用应前后带括号，通常使用 \cs{eqref} 命令，比如式\eqref{eq:example}。
\begin{equation}
  \frac{1}{2 \uppi \symup{i}} \int_\gamma f = \sum_{k=1}^m n(\gamma; a_k) \mathscr{R}(f; a_k).
  \label{eq:example}
\end{equation}

多行公式尽可能在“=”处对齐，推荐使用 \env{align} 环境。
\begin{align}
  a & = b + c + d + e \\
    & = f + g
\end{align}



\section{数学定理}

定理环境的格式可以使用 \pkg{amsthm} 或者 \pkg{ntheorem} 宏包配置。
用户在导言区载入这两者之一后，模板会自动配置 \env{theorem}、\env{proof} 等环境。

\begin{theorem}[Lindeberg--Lévy 中心极限定理]
  设随机变量 $X_1, X_2, \dots, X_n$ 独立同分布， 且具有期望 $\mu$ 和有限的方差 $\sigma^2 \ne 0$，
  记 $\bar{X}_n = \frac{1}{n} \sum_{i+1}^n X_i$，则
  \begin{equation}
    \lim_{n \to \infty} P \left(\frac{\sqrt{n} \left( \bar{X}_n - \mu \right)}{\sigma} \le z \right) = \Phi(z),
  \end{equation}
  其中 $\Phi(z)$ 是标准正态分布的分布函数。
\end{theorem}
\begin{proof}
  Trivial.
\end{proof}

同时模板还提供了 \env{assumption}、\env{definition}、\env{proposition}、
\env{lemma}、\env{theorem}、\env{axiom}、\env{corollary}、\env{exercise}、
\env{example}、\env{remar}、\env{problem}、\env{conjecture} 这些相关的环境。

% !TeX root = ../thuthesis-example.tex

\chapter{装药设计相关知识}

\section{推进剂选择原则}

推进剂应具有所需要的能量特性以比冲与密度的乘积来表示的。
\[
    I=I_{sp}m_pg=I_{sp}\rho _gg\mathrm{V}_p
\]
$V_{p}$一定时，$I_{sp}\rho _g$越大，$I$越大。相反$I$一定时，$I_{sp}\rho _g$越大，$V_{p}$越小，燃烧室体积也越小，燃烧室壳体质量也就越小。$I_{sp}$越大，$m_{p}$越小，发动机总质量也越小。所以，我们要求推进剂应该具有尽量高的能量特性。

\section{药形种类}

目前在发动机中采用的装药有很多，可以按照燃烧面积的变化规律、燃烧表面所处的位置划分，按侧面燃耗药柱划分，按燃烧方向划分。

按照燃烧面积的变化规律划分：
\begin{itemize}[leftmargin=2em]
    \item 恒面性药柱增面性药柱
    \item 减面性药柱
    \item 增面性药柱
\end{itemize}

按照燃烧表面所处的位置划分：
\begin{itemize}[leftmargin=2em]
    \item 端面燃烧药柱
    \item 侧面燃烧药柱
    \item 侧端面同时燃烧药柱
\end{itemize}

按照侧面燃耗药柱划分：
\begin{itemize}[leftmargin=2em]
    \item 外燃药柱
    \item 内燃药柱
    \item 内外燃药柱
\end{itemize}

按照燃烧方向划分：
\begin{itemize}[leftmargin=2em]
    \item 一维药柱
    \item 二维药柱
    \item 三维药柱
\end{itemize}

\subsection{一维药柱}

一维药柱即端燃药柱，这种药柱的侧表面及其一端是用包覆层阻燃的。燃烧只在一端进行，燃烧方向垂直于端面，因此是恒面性燃烧的药柱。

一维药柱具有恒面燃烧、工作时间长、装填系数最大、无压力峰、形状简单、制造容易、强度高等优点。但同时也具有推理小、重心偏移大、点火困难等不可避免的缺点。

\subsection{二维药柱}

二维药柱即侧燃药柱，由于侧燃药柱燃烧面积大，能够产生足够大的推力，在固体火箭发动机中被广泛采用。二维药柱种类很多，每种药柱都具有特定的燃烧面变化规律。

\begin{longtable}{c l l}
    \caption{不同二维药柱优缺点比较}
    \label{tab:longtable} \\
    \toprule
    种类 & \multicolumn{1}{c}{优点} & \multicolumn{1}{c}{缺点}  \\
    \midrule
  \endfirsthead
    \caption*{续表~\thetable\quad 不同二维药柱优缺点比较} \\
    \toprule
    种类 & \multicolumn{1}{c}{优点} & \multicolumn{1}{c}{缺点}  \\
    \midrule
  \endhead
    \bottomrule
  \endfoot
  
  \multirow{4}*{管型药柱} & 燃面恒定 & ~\\
 

  ~ & 结构简单，易于制造 & 药柱需要支撑装置 \\



  ~ & 无应力集中，强度大 &不可长时间工作 \\


  ~ & 无剩药 & \\

  \hline

  管套形药柱 & \begin{tabular}[c]{@{}l@{}}燃面恒定\\无剩药\end{tabular} & 支撑装置的工作条件恶劣\\
  \hline
  星形药柱 & \begin{tabular}[c]{@{}l@{}}可获得恒面性、减面性或增面性的燃烧\\工作时间可以很长\\可直接浇注在燃烧室内\\对壳体的刚度有增强作用\end{tabular} & \begin{tabular}[c]{@{}l@{}}药形复杂，药模制造困难\\药柱强度低，易出现裂纹\\有剩药\\推力、压力曲线\\有较长的拖尾现象\end{tabular}  \\
  \hline
  车轮形药柱 & \begin{tabular}[c]{@{}l@{}}可获得恒面性、减面性或增面性的燃烧\\工作时间可以很长\\肉厚系数低\\可以得到多推力方案\end{tabular}                                     & \begin{tabular}[c]{@{}l@{}}药柱强度低\\装填系数低，形状更为复杂\end{tabular}                              \\
  \hline
  多孔药柱  & 药柱绝热，可采用直接浇注工艺&\\
\end{longtable}

  


% 其他部分
\backmatter

% 参考文献
\bibliography{ref/refs}  % 参考文献使用 BibTeX 编译
% \printbibliography       % 参考文献使用 BibLaTeX 编译

% 附录
% 本科生需要将附录放到声明之后，个人简历之前
% \appendix
% % !TeX root = ../thuthesis-example.tex

\begin{survey}
\label{cha:survey}

\title{Title of the Survey}
\maketitle


\tableofcontents


本科生的外文资料调研阅读报告。


\section{Figures and Tables}

\subsection{Figures}

An example figure in appendix (Figure~\ref{fig:appendix-survey-figure}).

\begin{figure}
  \centering
  \includegraphics[width=0.6\linewidth]{example-image-a.pdf}
  \caption{Example figure in appendix}
  \label{fig:appendix-survey-figure}
\end{figure}


\subsection{Tables}

An example table in appendix (Table~\ref{tab:appendix-survey-table}).

\begin{table}
  \centering
  \caption{Example table in appendix}
  \begin{tabular}{ll}
    \toprule
    File name       & Description                                         \\
    \midrule
    thuthesis.dtx   & The source file including documentation and comments \\
    thuthesis.cls   & The template file                                   \\
    thuthesis-*.bst & BibTeX styles                                       \\
    thuthesis-*.bbx & BibLaTeX styles for bibliographies                  \\
    thuthesis-*.cbx & BibLaTeX styles for citations                       \\
    \bottomrule
  \end{tabular}
  \label{tab:appendix-survey-table}
\end{table}


\section{Equations}

An example equation in appendix (Equation~\eqref{eq:appendix-survey-equation}).
\begin{equation}
  \frac{1}{2 \uppi \symup{i}} \int_\gamma f = \sum_{k=1}^m n(\gamma; a_k) \mathscr{R}(f; a_k)
  \label{eq:appendix-survey-equation}
\end{equation}


\section{Citations}

Example\cite{dupont1974bone} citations\cite{merkt1995rotational} in appendix
\cite{dupont1974bone,merkt1995rotational}.


% 默认使用正文的参考文献样式；
% 如果使用 BibTeX，可以切换为其他兼容 natbib 的 BibTeX 样式。
\bibliographystyle{unsrtnat}
% \bibliographystyle{IEEEtranN}

% 默认使用正文的参考文献 .bib 数据库；
% 如果使用 BibTeX，可以改为指定数据库，如 \bibliography{ref/refs}。
\printbibliography

\end{survey}
       % 本科生：外文资料的调研阅读报告
% % !TeX root = ../thuthesis-example.tex

\begin{translation}
\label{cha:translation}

\title{书面翻译题目}
\maketitle

\tableofcontents


本科生的外文资料书面翻译。


\section{图表示例}

\subsection{图}

附录中的图片示例（图~\ref{fig:appendix-translation-figure}）。

\begin{figure}
  \centering
  \includegraphics[width=0.6\linewidth]{example-image-a.pdf}
  \caption{附录中的图片示例}
  \label{fig:appendix-translation-figure}
\end{figure}


\subsection{表格}

附录中的表格示例（表~\ref{tab:appendix-translation-table}）。

\begin{table}
  \centering
  \caption{附录中的表格示例}
  \begin{tabular}{ll}
    \toprule
    文件名          & 描述                         \\
    \midrule
    thuthesis.dtx   & 模板的源文件，包括文档和注释 \\
    thuthesis.cls   & 模板文件                     \\
    thuthesis-*.bst & BibTeX 参考文献表样式文件    \\
    thuthesis-*.bbx & BibLaTeX 参考文献表样式文件  \\
    thuthesis-*.cbx & BibLaTeX 引用样式文件        \\
    \bottomrule
  \end{tabular}
  \label{tab:appendix-translation-table}
\end{table}


\section{数学公式}

附录中的数学公式示例（公式\eqref{eq:appendix-translation-equation}）。
\begin{equation}
  \frac{1}{2 \uppi \symup{i}} \int_\gamma f = \sum_{k=1}^m n(\gamma; a_k) \mathscr{R}(f; a_k)
  \label{eq:appendix-translation-equation}
\end{equation}


\section{文献引用}

附录\cite{dupont1974bone}中的参考文献引用\cite{merkt1995rotational}示例
\cite{dupont1974bone,merkt1995rotational}。


\appendix

\section{附录}

附录的内容。


% 书面翻译的参考文献
% 默认使用正文的参考文献样式；
% 如果使用 BibTeX，可以切换为其他兼容 natbib 的 BibTeX 样式。
\bibliographystyle{unsrtnat}
% \bibliographystyle{IEEEtranN}

% 默认使用正文的参考文献 .bib 数据库；
% 如果使用 BibTeX，可以改为指定数据库，如 \bibliography{ref/refs}。
\printbibliography

% 书面翻译对应的原文索引
\begin{translation-index}
  \nocite{mellinger1996laser}
  \nocite{bixon1996dynamics}
  \nocite{carlson1981two}
  \bibliographystyle{unsrtnat}
  \printbibliography
\end{translation-index}

\end{translation}
  % 本科生：外文资料的书面翻译
% % !TeX root = ../thuthesis-example.tex

\chapter{补充内容}

附录是与论文内容密切相关、但编入正文又影响整篇论文编排的条理和逻辑性的资料，例如某些重要的数据表格、计算程序、统计表等，是论文主体的补充内容，可根据需要设置。

附录中的图、表、数学表达式、参考文献等另行编序号，与正文分开，一律用阿拉伯数字编码，
但在数码前冠以附录的序号，例如“图~\ref{fig:appendix-figure}”，
“表~\ref{tab:appendix-table}”，“式\eqref{eq:appendix-equation}”等。


\section{插图}

% 附录中的插图示例（图~\ref{fig:appendix-figure}）。

\begin{figure}
  \centering
  \includegraphics[width=0.6\linewidth]{example-image-a.pdf}
  \caption{附录中的图片示例}
  \label{fig:appendix-figure}
\end{figure}


\section{表格}

% 附录中的表格示例（表~\ref{tab:appendix-table}）。

\begin{table}
  \centering
  \caption{附录中的表格示例}
  \begin{tabular}{ll}
    \toprule
    文件名          & 描述                         \\
    \midrule
    thuthesis.dtx   & 模板的源文件，包括文档和注释 \\
    thuthesis.cls   & 模板文件                     \\
    thuthesis-*.bst & BibTeX 参考文献表样式文件    \\
    thuthesis-*.bbx & BibLaTeX 参考文献表样式文件  \\
    thuthesis-*.cbx & BibLaTeX 引用样式文件        \\
    \bottomrule
  \end{tabular}
  \label{tab:appendix-table}
\end{table}


\section{数学表达式}

% 附录中的数学表达式示例（式\eqref{eq:appendix-equation}）。
\begin{equation}
  \frac{1}{2 \uppi \symup{i}} \int_\gamma f = \sum_{k=1}^m n(\gamma; a_k) \mathscr{R}(f; a_k)
  \label{eq:appendix-equation}
\end{equation}


\section{文献引用}

附录\cite{dupont1974bone}中的参考文献引用\cite{zhengkaiqing1987}示例
\cite{dupont1974bone,zhengkaiqing1987}。

\printbibliography


% 致谢
% !TeX root = ../thuthesis-example.tex

\begin{acknowledgements}
  感谢帮助我解决问题的挚友。
\end{acknowledgements}


% 声明
% 本科生开题报告不需要
% \statement
% 将签字扫描后的声明文件 scan-statement.pdf 替换原始页面
% \statement[file=scan-statement.pdf]
% 本科生编译生成的声明页默认不加页脚，插入扫描版时再补上；
% 研究生编译生成时有页眉页脚，插入扫描版时不再重复。
% 也可以手动控制是否加页眉页脚
% \statement[page-style=empty]
% \statement[file=scan-statement.pdf, page-style=plain]

% % 个人简历、在学期间完成的相关学术成果
% % 本科生可以附个人简历，也可以不附个人简历
% % !TeX root = ../thuthesis-example.tex

\begin{resume}

  \section*{个人简历}

  197× 年 ×× 月 ×× 日出生于四川××县。

  1992 年 9 月考入××大学化学系××化学专业，1996 年 7 月本科毕业并获得理学学士学位。

  1996 年 9 月免试进入清华大学化学系攻读××化学博士至今。


  \section*{在学期间完成的相关学术成果}

  \subsection*{学术论文}

  \begin{achievements}
    \item Yang Y, Ren T L, Zhang L T, et al. Miniature microphone with silicon-based ferroelectric thin films[J]. Integrated Ferroelectrics, 2003, 52:229-235.
    \item 杨轶, 张宁欣, 任天令, 等. 硅基铁电微声学器件中薄膜残余应力的研究[J]. 中国机械工程, 2005, 16(14):1289-1291.
    \item 杨轶, 张宁欣, 任天令, 等. 集成铁电器件中的关键工艺研究[J]. 仪器仪表学报, 2003, 24(S4):192-193.
    \item Yang Y, Ren T L, Zhu Y P, et al. PMUTs for handwriting recognition. In press[J]. (已被Integrated Ferroelectrics录用)
  \end{achievements}


  \subsection*{专利}

  \begin{achievements}
    \item 任天令, 杨轶, 朱一平, 等. 硅基铁电微声学传感器畴极化区域控制和电极连接的方法: 中国, CN1602118A[P]. 2005-03-30.
    \item Ren T L, Yang Y, Zhu Y P, et al. Piezoelectric micro acoustic sensor based on ferroelectric materials: USA, No.11/215, 102[P]. (美国发明专利申请号.)
  \end{achievements}

\end{resume}


% % 指导教师/指导小组评语
% % 本科生不需要
% % !TeX root = ../thuthesis-example.tex

\begin{comments}
% \begin{comments}[name = {指导小组评语}]
% \begin{comments}[name = {Comments from Thesis Supervisor}]
% \begin{comments}[name = {Comments from Thesis Supervision Committee}]

  论文提出了……

\end{comments}


% % 答辩委员会决议书
% % 本科生不需要
% % !TeX root = ../thuthesis-example.tex

\begin{resolution}

  论文提出了……

  论文取得的主要创新性成果包括：

  1. ……

  2. ……

  3. ……

  论文工作表明作者在×××××具有×××××知识，具有××××能力，论文××××，答辩××××。

  答辩委员会表决，（×票/一致）同意通过论文答辩，并建议授予×××（姓名）×××（门类）学博士/硕士学位。

\end{resolution}


% 本科生的综合论文训练记录表（扫描版）
% \record{file=scan-record.pdf}

\end{document}
